\unit{Bài trình bày 195 -- Kỹ thuật hạ xuống II}
\sid{195}
\markboth{Thuyết trình 195}{Thuyết trình 195}
\src{369}
\seminar{Năm thứ 12, 1959/60, số 195\\Tháng 2 năm 1960}
\paperhead{KỸ THUẬT HẠ XUỐNG VÀ CÁC ĐỊNH LÝ TỒN TẠI TRONG HÌNH HỌC ĐẠI SỐ\\
II. ĐỊNH LÝ TỒN TẠI TRONG LÝ THUYẾT HÌNH THỨC VỀ MÔĐUN}%
  {do Alexander GROTHENDIECK}

\fsection{A}{Các hàm tử biểu diễn được và pro-biểu diễn được}

\subsection*{1. Các hàm tử biểu diễn được}
\addcontentsline{toc}{subsection}{1. Các hàm tử biểu diễn được}

Cho $\mathcal C$ là một phạm trù. Với mọi $X\in\mathcal C$, ký hiệu $h_X$ là
hàm tử phản biến từ $\mathcal C$ đến phạm trù $(\operatorname{Ens})$ của các
tập hợp,
\[
  h_X:\mathcal C\longrightarrow(\operatorname{Ens}),
\]
được xác định bởi công thức
\[
  h_X(Y)=\operatorname{Hom}(Y,X).
\]
Một cấu xạ $X\longrightarrow X'$ trong $\mathcal C$ hiển nhiên cảm sinh một
đồng cấu hàm tử $h_X\longrightarrow h_{X'}$: $h_X$ là một hàm tử hiệp biến
theo $X$; nói cách khác, ta đã xác định một hàm tử hiệp biến chính tắc
\[
  h:\mathcal C\longrightarrow
  \operatorname{Hom}\bigl(\mathcal C^o,(\operatorname{Ens})\bigr)
\]
từ $\mathcal C$ đến phạm trù các hàm tử hiệp biến từ phạm trù đối ngẫu
$\mathcal C^o$ của $\mathcal C$ đến phạm trù các tập hợp. Ta nhắc lại:

\result{MỆNH ĐỀ}{1.1} Hàm tử $h$ này \emph{đầy đủ và trung thành}; nói cách
khác, với mọi cặp đối tượng $X,X'$ của $\mathcal C$, ánh xạ tự nhiên
\[
  \operatorname{Hom}(X,X')\longrightarrow
  \operatorname{Hom}(h_X,h_{X'})
\]
là song ánh.

Đặc biệt, nếu một hàm tử
$F\in\operatorname{Hom}(\mathcal C^o,(\operatorname{Ens}))$ đẳng cấu với một
hàm tử dạng $h_X$, thì $X$ được xác định sai khác một đẳng cấu duy nhất.
Trong trường hợp đó ta nói hàm tử $F$ là \emph{biểu diễn được}. Mệnh đề trên
có nghĩa rằng hàm tử chính tắc $h$ xác lập một tương đương giữa phạm trù
$\mathcal C$ và phạm trù con đầy đủ của
$\operatorname{Hom}(\mathcal C^o,(\operatorname{Ens}))$ gồm các hàm tử biểu
diễn được. Sự kiện này là nền tảng của khái niệm «nghiệm của một bài toán phổ
quát»: một bài toán như vậy luôn quy về việc xét xem một hàm tử đã cho
(phản biến như ở đây, hoặc hiệp biến trong trường hợp đối ngẫu) từ
$\mathcal C$ đến $(\operatorname{Ens})$ có biểu diễn được hay không.

\src{370}
Ngoài ra, theo chính định nghĩa của tích trong một phạm trù [1], hàm tử
$h:X\longmapsto h_X$ giao hoán với các tích mỗi khi chúng tồn tại (và tổng
quát hơn, với các giới hạn xạ ảnh hữu hạn hoặc vô hạn; đặc biệt với các tích
thớ, với phép lập «hạt nhân» [2], v.v., mỗi khi chúng tồn tại): ta có một
đẳng cấu hàm tử
\[
  h_{X\times X'}\xrightarrow{\sim}h_X\times h_{X'}
\]
mỗi khi $X\times X'$ tồn tại; nghĩa là ta có các song ánh hàm tử theo $Y$
\[
  h_{X\times X'}(Y)\xrightarrow{\sim}h_X(Y)\times h_{X'}(Y).
\]
Đặc biệt, việc cho một cấu xạ
\[
  X\times X'\longrightarrow X''
\]
trong $\mathcal C$ (tức một «luật hợp thành» trong $\mathcal C$ giữa
$X,X',X''$) tương đương với việc cho một cấu xạ
\[
  h_{X\times X'}=h_X\times h_{X'}\longrightarrow h_{X''},
\]
tức là, với mọi $Y\in\mathcal C$, cho một luật hợp thành các tập hợp
\[
  h_X(Y)\times h_{X'}(Y)\longrightarrow h_{X''}(Y)
\]
sao cho, với mọi cấu xạ $Y\longrightarrow Y'$ trong $\mathcal C$, hệ các ánh
xạ giữa các tập hợp
\[
  h_{X^{(i)}}(Y)\longleftarrow h_{X^{(i)}}(Y')
  \qquad(\text{với }i=0,1,2)
\]
là một đồng cấu đối với hai luật hợp thành tương ứng với $Y$ và $Y'$. Do đó,
khái niệm cấu trúc «$\mathcal C$-nhóm», «$\mathcal C$-vành», v.v. trên một
đối tượng $X$ của $\mathcal C$ được diễn đạt thuận tiện nhất (cả trong lý
thuyết lẫn thực hành) bằng cách nói rằng, với mọi $Y\in\mathcal C$, tập
$h_X(Y)$ được trang bị một luật nhóm (tương ứng, luật vành, v.v.) theo nghĩa
thông thường, và các ánh xạ
$h_X(Y)\longleftarrow h_X(Y')$ tương ứng với các cấu xạ
$Y\longrightarrow Y'$ là các đồng cấu nhóm (tương ứng, đồng cấu vành, v.v.).
Chẳng hạn, đây là cách trực quan và thuận tiện nhất để xác định các nhóm cổ
điển khác nhau $G_a$, $G_m$, $\operatorname{Gl}(n)$, v.v. trên một tiền lược
đồ cơ sở tùy ý $S$ và viết các quan hệ cổ điển giữa chúng; để trang bị cho
lược đồ afin $V(\mathcal F)$ trên $S$, được xác định bởi một bó tựa kết hợp
$\mathcal F$, một cấu trúc «phân thớ vectơ»; và để xác định rồi nghiên cứu
các đa tạp cờ liên hệ khác nhau (các grassmannian, các phân thớ xạ ảnh),
v.v. Nguyên lý chung chỉ đơn giản là đồng nhất, thông qua hàm tử chính tắc
$h$, các đối tượng của $\mathcal C$ với một số hàm tử phản biến: các hàm tử
biểu diễn được,

\src{371}
từ $\mathcal C$ đến phạm trù các tập hợp.

Thủ tục quen thuộc đảo chiều các mũi tên, cần thiết chẳng hạn trong trường hợp
các lược đồ afin để chuyển từ ngôn ngữ hình học sang ngôn ngữ đại số giao
hoán, dẫn đến việc đối ngẫu hóa các xét đoán trên và, đặc biệt, cũng đưa vào
các hàm tử hiệp biến biểu diễn được
$\mathcal C\longrightarrow(\operatorname{Ens})$, tức là các hàm tử dạng
\[
  Y\longmapsto\operatorname{Hom}(X,Y)=h'_X(Y).
\]

\subsection*{2. Các hàm tử pro-biểu diễn được, các pro-đối tượng}
\addcontentsline{toc}{subsection}{2. Các hàm tử pro-biểu diễn được, các pro-đối tượng}

Cho $\widetilde X=(X_i)_{i\in I}$ là một hệ xạ ảnh các đối tượng của
$\mathcal C$; hệ này tương ứng với một hàm tử hiệp biến
\[
  h'_{\widetilde X}=\varinjlim_i h'_{X_i},
\]
hay tường minh hơn,
\[
  h'_{\widetilde X}(Y)=\varinjlim_i h'_{X_i}(Y)
  =\varinjlim_i\operatorname{Hom}(X_i,Y),
\]
từ $\mathcal C$ đến $(\operatorname{Ens})$. Một hàm tử từ $\mathcal C$ đến
$(\operatorname{Ens})$ đẳng cấu với một hàm tử loại này, với $I$ lọc, được
gọi là \emph{pro-biểu diễn được}. Theo mục trước, đó chính là các hàm tử đẳng
cấu với giới hạn quy nạp lọc của các hàm tử biểu diễn được. Cho
$\widetilde X'=(X'_j)_{j\in J}$ là một hệ xạ ảnh lọc thứ hai trong
$\mathcal C$ (được xây dựng trên một tập chỉ số tiền thứ tự lọc thứ hai
$J$); dễ kiểm tra rằng có một song ánh chính tắc (tổng quát hóa Mệnh đề 1.1)
\[
  \operatorname{Hom}(h'_{\widetilde X'},h'_{\widetilde X})
  =\varprojlim_j\varinjlim_i\operatorname{Hom}(X_i,X'_j).
\]
Điều này dẫn ta đến phạm trù $\operatorname{Pro}(\mathcal C)$ của các
pro-đối tượng của $\mathcal C$; các đối tượng của nó là các hệ xạ ảnh các đối
tượng của $\mathcal C$ (trên các tập chỉ số tiền thứ tự lọc thay đổi), và nếu
$\widetilde X=(X_i)_{i\in I}$ cùng
$\widetilde X'=(X'_j)_{j\in J}$ là hai đối tượng thuộc lớp này, ta đặt
\[
  \operatorname{ProHom}(\widetilde X,\widetilde X')
  =\varprojlim_j\varinjlim_i\operatorname{Hom}(X_i,X'_j),
\]
với phép hợp thành hiển nhiên của các pro-đồng cấu. Theo chính cách xây dựng,
$\widetilde X\longmapsto h'_{\widetilde X}$ có thể được xem như một hàm tử
phản biến theo $\widetilde X$, xác lập một tương đương giữa phạm trù đối ngẫu
của phạm trù $\operatorname{Pro}(\mathcal C)$ các pro-đối tượng của
$\mathcal C$ và phạm trù các hàm tử hiệp biến pro-biểu diễn được từ
$\mathcal C$ đến $(\operatorname{Ens})$. Dĩ nhiên, một đối tượng $X$ của
$\mathcal C$ xác định chính tắc một pro-đối tượng, cũng được ký hiệu là $X$,

\src{372}
sao cho $\mathcal C$ tương đương với một phạm trù con đầy đủ của
$\operatorname{Pro}(\mathcal C)$. Khi đó, nếu
$\widetilde X=(X_i)_{i\in I}$ là một pro-đối tượng tùy ý của $\mathcal C$, ta
có (với sự đồng nhất trên)
\[
  \widetilde X=\varprojlim_i X_i,
\]
trong đó giới hạn xạ ảnh được lấy trong $\operatorname{Pro}(\mathcal C)$
(vì $h'_{\widetilde X}=\varinjlim_i h'_{X_i}$). Cần lưu ý rằng, ngay cả khi
giới hạn xạ ảnh của các $X_i$ tồn tại trong $\mathcal C$, nói chung nó không
đẳng cấu với giới hạn xạ ảnh $\widetilde X$ trong
$\operatorname{Pro}(\mathcal C)$, như đã hiển nhiên khi $\mathcal C$ là phạm
trù các tập hợp. Theo định nghĩa,
$\varprojlim_{\mathcal C}X_i=L$ được xác định bởi điều kiện rằng hàm tử theo
$Y\in\mathcal C$
\[
  \varprojlim_i\operatorname{Hom}_{\mathcal C}(Y,X_i)
  =\operatorname{Hom}_{\operatorname{Pro}(\mathcal C)}(Y,\widetilde X)
\]
có giá trị trong $(\operatorname{Ens})$ được biểu diễn bởi $L$, tức là đẳng
cấu với $\operatorname{Hom}_{\mathcal C}(Y,L)$; do đó,
$\varprojlim_{\mathcal C}X_i$ đã được xác định theo pro-đối tượng
$\widetilde X$, và chính xác hơn, phụ thuộc hàm tử vào pro-đối tượng
$\widetilde X$ mỗi khi nó được xác định; vì vậy không có trở ngại gì khi ký
hiệu nó là $\varprojlim_{\mathcal C}(\widetilde X)$. Nếu các giới hạn xạ ảnh
luôn tồn tại trong $\mathcal C$,
$\varprojlim_{\mathcal C}(\widetilde X)$ là một hàm tử từ
$\operatorname{Pro}(\mathcal C)$ đến $\mathcal C$, và có một đồng cấu hàm tử
chính tắc
$\varprojlim_{\mathcal C}(\widetilde X)\longrightarrow\widetilde X$.

Mọi hàm tử, mà để định ý ta giả sử là hiệp biến,
\[
  F:\mathcal C\longrightarrow\mathcal C'
\]
đều được mở rộng một cách hiển nhiên thành hàm tử
\[
  \operatorname{Pro}(F):\operatorname{Pro}(\mathcal C)
  \longrightarrow\operatorname{Pro}(\mathcal C'),
\]
nên nếu các giới hạn xạ ảnh luôn tồn tại trong $\mathcal C'$, thì $F$ cũng
xác định chính tắc một hàm tử hợp thành
\[
  \overline F:\operatorname{Pro}(\mathcal C)\longrightarrow\mathcal C',
  \qquad
  \overline F=\varprojlim_{\mathcal C'}\circ\operatorname{Pro}(F),
\]
biến $\widetilde X=(X_i)_{i\in I}$ thành
$\varprojlim_{\mathcal C'}F(X_i)$.

Một pro-đối tượng $\widetilde X$ được gọi là \emph{pro-đối tượng nghiêm ngặt}
nếu nó đẳng cấu với một pro-đối tượng $(X_i)_{i\in I}$ mà các cấu xạ chuyển
$X_i\longrightarrow X_j$ đều là toàn cấu; một hàm tử được xác định bởi một
đối tượng như vậy được gọi là \emph{pro-biểu diễn được nghiêm ngặt}. Ngoài
ra, có thể đòi hỏi $I$ là tập có thứ tự lọc và mọi toàn cấu
$X_i\longrightarrow X'$ tương đương với một toàn cấu

\src{373}
$X_i\longrightarrow X_j$ với một $j\in I$ thích hợp (được điều kiện này xác
định duy nhất). Dưới các điều kiện ấy, hệ xạ ảnh $(X_i)_{i\in I}$ được xác
định sai khác một đẳng cấu duy nhất (theo nghĩa thông thường của đẳng cấu các
hệ xạ ảnh). Suy ra rằng trong $\operatorname{Pro}(\mathcal C)$ luôn tồn tại
giới hạn xạ ảnh của một hệ xạ ảnh các pro-đối tượng nghiêm ngặt
$\widetilde X^{(\alpha)}$, và với các ký hiệu trên cho $F$, $\overline F$, ta
có
\[
  \overline F\left(\varprojlim_\alpha\widetilde X^{(\alpha)}\right)
  =\varprojlim_{\alpha,\mathcal C'}
  \overline F\left(\widetilde X^{(\alpha)}\right).
\]
Đặc biệt, nếu mọi pro-đối tượng của $\mathcal C$ đều nghiêm ngặt (xem mục
sau), thì hàm tử mở rộng $\overline F$ giao hoán với các giới hạn xạ ảnh.

\subsection*{3. Đặc trưng các hàm tử pro-biểu diễn được}
\addcontentsline{toc}{subsection}{3. Đặc trưng các hàm tử pro-biểu diễn được}

Cho $\mathcal C$, $\mathcal C'$ là hai phạm trù trong đó các giới hạn xạ ảnh
hữu hạn tồn tại (tức là các giới hạn tương ứng với các tập tiền thứ tự hữu
hạn, không nhất thiết lọc), hay tương đương, trong đó các tích hữu hạn và các
tích thớ hữu hạn tồn tại (điều này đặc biệt kéo theo sự tồn tại của một «đối
tượng cuối» $e$, sao cho $\operatorname{Hom}(X,e)$ chỉ gồm một phần tử với
mọi $X$). Cho $F$ là một hàm tử hiệp biến từ $\mathcal C$ đến
$\mathcal C'$. Các điều kiện sau tương đương:

\begin{enumerate}
\item[(i)] $F$ giao hoán với các giới hạn xạ ảnh hữu hạn;
\item[(ii)] $F$ giao hoán với các tích hữu hạn và các tích thớ hữu hạn;
\item[(iii)] $F$ giao hoán với các tích hữu hạn và, với mọi biểu đồ khớp
\[
  X\longrightarrow X'\rightrightarrows X''
\]
trong $\mathcal C$ ([3], A, Định nghĩa 2.1), biểu đồ được biến đổi bởi $F$
\[
  F(X)\longrightarrow F(X')\rightrightarrows F(X'')
\]
là khớp.
\end{enumerate}

Khi đó ta nói $F$ \emph{khớp trái}.

Từ đây trở đi, ta giả sử các giới hạn xạ ảnh hữu hạn tồn tại trong
$\mathcal C$. Từ các định nghĩa, lập tức suy ra rằng một hàm tử biểu diễn
được là khớp trái và, khi chuyển sang giới hạn, một hàm tử pro-biểu diễn được
là khớp trái.

Để có một mệnh đề đảo, xét
\[
  F:\mathcal C\longrightarrow(\operatorname{Ens})
\]
là một hàm tử hiệp biến, và cho $X\in\mathcal C$, $\xi\in F(X)$. Ta nói
$\xi$ (hay cặp \src{374}$(X,\xi)$) là \emph{tối tiểu} nếu, với mọi
$X'\in\mathcal C$, $\xi'\in F(X')$ và mọi đơn cấu nghiêm ngặt ([3], A, 2)
$u:X'\longrightarrow X$ sao cho $\xi=F(u)(\xi')$, thì $u$ là một đẳng cấu.
Mặt khác, ta nói một cặp $(X,\xi)$ \emph{trội} $(X'',\xi'')$
($\xi\in F(X)$, $\xi''\in F(X'')$) nếu tồn tại một cấu xạ
$v:X\longrightarrow X''$ sao cho $\xi''=F(v)(\xi)$; nếu $\xi$ tối tiểu và
$F$ khớp trái, thì cấu xạ $v$ là duy nhất; nếu $\xi''$ tối tiểu, thì $v$ là
toàn ánh. Dễ dàng suy ra mệnh đề sau:

\result{MỆNH ĐỀ}{3.1} Để $F$ pro-biểu diễn được nghiêm ngặt, điều kiện cần và
đủ là nó thỏa mãn hai điều kiện:
\begin{enumerate}
\item[(i)] $F$ khớp trái;
\item[(ii)] mọi cặp $(X,\xi)$, với $\xi\in F(X)$, đều bị trội bởi một cặp
tối tiểu.
\end{enumerate}

Điều kiện cuối là trống nếu mọi đối tượng của $\mathcal C$ đều là Artin (lấy
một đối tượng con $X'$ của $X$ tối tiểu trong số những đối tượng mà với chúng
tồn tại $\xi'\in F(X')$ sao cho $\xi$ là ảnh của $\xi'$). Do đó:

\result{HỆ QUẢ}{} Cho $\mathcal C$ là một phạm trù mà các đối tượng đều là
Artin và trong đó các giới hạn xạ ảnh hữu hạn tồn tại. Khi đó, các hàm tử
pro-biểu diễn được từ $\mathcal C$ đến $(\operatorname{Ens})$ chính là các
hàm tử khớp trái, và thậm chí đều pro-biểu diễn được nghiêm ngặt.

Sự kiện cuối này cũng có nghĩa rằng \emph{mọi pro-đối tượng của $\mathcal C$
đều nghiêm ngặt}.

\subsection*{4. Ví dụ: các nhóm kiểu Galois, các nhóm pro-đại số}
\addcontentsline{toc}{subsection}{4. Các nhóm kiểu Galois, các nhóm pro-đại số}

Nếu $\mathcal C$ là phạm trù các nhóm hữu hạn thông thường, thì
$\operatorname{Pro}(\mathcal C)$ tương đương với phạm trù các nhóm tôpô
compact hoàn toàn không liên thông. Các nhóm loại này và các tổng quát hóa
của chúng, thu được bằng cách thay các nhóm hữu hạn thông thường bởi các lược
đồ nhóm hữu hạn trên một tiền lược đồ cơ sở đã cho (chẳng hạn, các nhóm đại
số hữu hạn trên một trường $k$), sẽ đóng vai trò các nhóm cơ bản, đồng luân
và đồng điều tuyệt đối và tương đối của các tiền lược đồ. Trong tất cả các
ví dụ này, Hệ quả của Mệnh đề 3.1 được áp dụng, và chính thông qua hàm tử liên
hệ mà ta phải xác định $\pi_1$ [2]. Điều tương tự xảy ra khi bắt đầu từ phạm
trù các nhóm đại số hoặc tựa đại số trên một trường (hay tổng quát hơn, trên
một tiền lược đồ Noether): ta thu được các «nhóm pro-đại số» của Serre [4].

\subsection*{5. Ví dụ: «các đa tạp hình thức»}
\addcontentsline{toc}{subsection}{5. Ví dụ: các đa tạp hình thức}

Cho $\Lambda$ là một vành Noether và $\mathcal C$ là phạm trù các
$\Lambda$-đại số $A$ là những môđun Artin kiểu hữu hạn trên $\Lambda$ (hay
gọi tắt là

\src{375}
các $\Lambda$-đại số Artin). Các điều kiện của Hệ quả thuộc Mệnh đề 3.1 được
thỏa mãn. Phạm trù $\operatorname{Pro}(\mathcal C)$ ở đây tương đương với
phạm trù các đại số tôpô $O$ trên $\Lambda$ đẳng cấu với các giới hạn xạ ảnh
tôpô
\[
  O=\varprojlim_i O_i
\]
của các đại số $O_i\in\mathcal C$, tức là tôpô của chúng tuyến tính, tách và
đầy đủ, đồng thời với mọi iđêan mở $J_i$ của $O$, $O/J_i$ là một đại số Artin
trên $\Lambda$. Hàm tử
$\mathcal C\longrightarrow(\operatorname{Ens})$ liên hệ với một đại số như
vậy chính là
\begin{align*}
  F(A)=h'_O(A)
  &=\{\text{các đồng cấu liên tục của các $\Lambda$-đại số tôpô }
       O\longrightarrow A\}\\
  &=\varinjlim_i
    \operatorname{Hom}_{\Lambda\text{-đại số}}(O_i,A).
\end{align*}

Ngoài ra, lưu ý rằng phạm trù $\mathcal C$ về cơ bản là tích của các phạm trù
tương tự ứng với các phép đầy đủ hóa của những vành địa phương
$\Lambda_{\mathfrak m}$, với $\mathfrak m$ chạy qua các iđêan tối đại của
$\Lambda$; do đó, nếu muốn, có thể giới hạn vào trường hợp $\Lambda$ là một
vành địa phương đầy đủ thuộc lớp này. Trong mọi trường hợp, $O$ phân rã chính
tắc thành tích tôpô của các thành phần địa phương, tương ứng với các «điểm»
của lược đồ hình thức [2] được xác định bởi $O$. Một điểm thuộc lớp này được
xác định bởi một phần tử $\xi$ của một $F(K)$ nào đó, trong đó
$K\in\mathcal C$ là một trường (chẳng hạn trường thặng dư của thành phần địa
phương đang xét); hai cặp $(\xi,K)$ và $(\xi',K')$ xác định cùng một điểm khi
và chỉ khi chúng bị trội bởi cùng một $(\xi'',K'')$, hoặc cùng trội một
$(\xi''',K''')$. (Nếu các trường $\Lambda/\mathfrak m$ đóng đại số, chỉ cần
lấy hợp rời của các $F(\Lambda/\mathfrak m)$.)

Điều quan trọng là cho các điều kiện để thành phần địa phương $O_\xi$ của
$O$ ứng với $\xi\in F(K)$ là một vành Noether. Khi $\Lambda$ là một vành địa
phương đầy đủ (nhắc lại là Noether), điều này tương đương với việc $O_\xi$
đẳng cấu với một vành thương của vành chuỗi lũy thừa hình thức
$\Lambda[[t_1,\ldots,t_n]]$ trên $\Lambda$. Để nêu một tiêu chuẩn như vậy,
với mỗi vành $A$ ta đưa vào $A$-đại số $I_A$ các «số đối ngẫu» của $A$ bằng
công thức
\[
  I_A=A[t]/t^2A[t].
\]
Cho $\varepsilon:I_A\longrightarrow A$ là đồng cấu tăng cường; đồng cấu này
xác định (nếu $A\in\mathcal C$) một ánh xạ
\[
  F(\varepsilon):F(I_A)\longrightarrow F(A),
\]

\src{376}
và, bằng cách sử dụng tính khớp trái của $F$, ta xác định một cách nội tại
một cấu trúc $K$-môđun trên phần
\[
  F(I_A,\xi)=F(\varepsilon)^{-1}(\xi)
\]
của $F(I_A)$ gồm các $\xi'\in F(I_A)$ «rút gọn thành $\xi$»; từ dạng tường
minh của $F$ theo $O$, ta kiểm tra được rằng $A$-môđun này đồng nhất với
\[
  \operatorname{Hom}_{\Lambda}(\mathfrak m_\xi/\mathfrak m_\xi^2,A),
\]
trong đó $\mathfrak m_\xi$ là hạt nhân của đồng cấu
$\xi:O\longrightarrow A$, tức là, nếu $A$ là một trường, iđêan tối đại của
thành phần địa phương $O_\xi$ của $O$. Ta lập tức suy ra mệnh đề sau:

\result{MỆNH ĐỀ}{5.1} Cho $\xi\in F(K)$, trong đó $K\in\mathcal C$ là một
trường. Để thành phần địa phương tương ứng $O_\xi$ của $O$ là một vành
Noether, điều kiện cần và đủ là tập $F(I_K,\xi)$ các phần tử của $F(I_K)$ rút
gọn thành $\xi$ là một không gian vectơ hữu hạn chiều trên $K$. Khi đó, ta có
một đẳng cấu chính tắc
\[
  F(I_K,\xi)=
  \operatorname{Hom}\!\left(
    \mathfrak m_\xi/\mathfrak m_\xi^2+\mathfrak n_\xi O_\xi,K
  \right),
\]
trong đó $\mathfrak n_\xi$ là iđêan tối đại của $\Lambda$, tức hạt nhân của
đồng cấu $\Lambda\longrightarrow K$.
Đặc biệt, chiều của $K$-không gian vectơ $F(I_K,\xi)$ bằng chiều của không
gian vectơ $\mathfrak m_\xi/\mathfrak m_\xi^2$ trên trường
$O_\xi/\mathfrak m_\xi=K(\xi)$.

Giả sử $O_\xi$ là Noether và, để đơn giản ký hiệu, $\Lambda$ là địa phương
đầy đủ và $O=O_\xi$. Ta nói $O$ \emph{trơn trên $\Lambda$} nếu $O$ là một
đại số hữu hạn và étale trên đại số đầy đủ hóa của phép địa phương hóa
$\Lambda[t_1,\ldots,t_n]$ tại một iđêan tối đại của vành này, iđêan ấy cảm
sinh iđêan tối đại của $\Lambda$; khi mở rộng thặng dư của $O$ trên
$\Lambda$ là tầm thường (chẳng hạn, nếu trường thặng dư của $\Lambda$ đóng
đại số), điều này tương đương với việc chính $O$ đẳng cấu với một đại số
chuỗi lũy thừa hình thức loại đó.
Cuối cùng, nếu không còn giả sử $O$ nhất thiết Noether, ta cũng nói $O$ trơn
trên $\Lambda$ nếu $O$ đẳng cấu với một giới hạn xạ ảnh tôpô của các
$\Lambda$-đại số thương Noether và trơn trên $\Lambda$ theo nghĩa trên. Điều
này lập tức tổng quát hóa cho trường hợp không giả sử $\Lambda$ và $O$ địa
phương. Khi đó ta có mệnh đề sau:

\result{MỆNH ĐỀ}{5.2} Để $O$ trơn trên $\Lambda$, điều kiện cần và đủ là hàm
tử liên hệ $F$ biến các toàn cấu thành các toàn cấu.

Điều này có nghĩa rằng với mọi đồng cấu toàn ánh
$A\longrightarrow A'$ trong $\mathcal C$, ánh xạ
$F(A)\longrightarrow F(A')$ là toàn ánh. Dĩ nhiên, chỉ cần kiểm tra điều kiện
này khi $A$ địa phương và, tiến hành từng bước, khi iđêan của $A$ là hạt nhân
của $A\longrightarrow A'$ bị triệt tiêu bởi iđêan tối đại của $A$. Trong
thực hành, điều này quy bài toán về việc kiểm tra rằng một cản trở nào đó,
gắn với các bất biến vô cùng bé

\src{377}
của tình huống sinh ra hàm tử $F$, bằng không; đây là một bài toán có bản
chất đối đồng điều.

Để kết thúc, ta nói đôi lời trong bối cảnh trên về các \emph{vành xác định}.
Vẫn cho $F$ là một hàm tử từ $\mathcal C$ đến $(\operatorname{Ens})$,
pro-biểu diễn được bởi một $\Lambda$-đại số tôpô $O$. Khi đó, với mọi
$A\in\mathcal C$ và mọi $\xi\in F(A)$, tồn tại một vành con nhỏ nhất $A'$ của
$A$ sao cho $\xi$ là ảnh của một phần tử $\xi'$ của $F(A')$ (khi đó phần tử
này được xác định duy nhất): thật vậy, chỉ cần diễn giải $\xi$ bằng một đồng
cấu từ $O$ đến $A$ rồi lấy $A'$ là ảnh của $O$ qua đồng cấu ấy. Khi đó ta
gọi $A'$ là vành xác định của đối tượng $\xi\in F(A)$. Nếu
$u:A\longrightarrow B$ là một đồng cấu đại số và $\eta=F(u)(\xi)$, thì vành
xác định của $\eta$ là ảnh qua $u$ của vành xác định của $\xi$. Khi bắt đầu
từ một hàm tử $F$ từ $\mathcal C$ đến $(\operatorname{Ens})$, sự tồn tại của
các vành xác định cùng các tính chất vừa nêu, với một vài dè dặt nhỏ, tương
đương với việc $F$ pro-biểu diễn được; nghĩa là, nói chung chúng còn lâu mới
tầm thường.

\fsection{B}{Hai định lý tồn tại}

Giữ các ký hiệu của A, mục 5, và xuất phát từ một hàm tử hiệp biến
\[
  F:\mathcal C\longrightarrow(\operatorname{Ens});
\]
ta tìm các tiêu chuẩn dễ sử dụng để $F$ pro-biểu diễn được, tức là có thể
diễn đạt bằng một $\Lambda$-đại số $O$ như trên. Theo Hệ quả của A, Mệnh đề
3.1, điều kiện cần và đủ là $F$ khớp trái. Ở tình trạng hiện nay của kỹ thuật
hạ xuống (xem các câu hỏi nêu trong [3], tr. 9), trong những trường hợp quan
trọng nhất không thể kiểm tra trực tiếp tiêu chuẩn này dưới dạng đó, và cần
các tiêu chuẩn có vẻ ít đòi hỏi hơn.

\result{ĐỊNH LÝ}{1} Để hàm tử $F$ pro-biểu diễn được, điều kiện cần và đủ là
nó thỏa mãn hai điều kiện sau:
\begin{enumerate}
\item[(i)] $F$ giao hoán với các tích hữu hạn;
\item[(ii)] với mọi đại số $A\in\mathcal C$ và mọi đồng cấu
$A\longrightarrow A'$ trong $\mathcal C$ sao cho biểu đồ
\[
  A\longrightarrow A'\rightrightarrows A'\otimes_A A'
\]
là khớp ([3], A, Định nghĩa 1.2), biểu đồ biến đổi
\[
  F(A)\longrightarrow F(A')\rightrightarrows F(A'\otimes_A A')
\]

\src{378}
là khớp.
\end{enumerate}
Hơn nữa, chỉ cần kiểm tra (ii) khi $A$ địa phương và thêm vào đó thuộc một
trong hai trường hợp sau:
\begin{enumerate}
\item[(a)] $A'$ là một môđun tự do trên $A$;
\item[(b)] môđun thương $A'/A$ là một $A$-môđun có độ dài 1.
\end{enumerate}

Chứng minh định lý này khá tinh tế và không thể phác thảo ở đây. Ta chỉ nêu
rằng nó chủ yếu dựa trên việc nghiên cứu các quan hệ tương đương (theo nghĩa
phạm trù) trên phổ của một đại số Artin; nghiên cứu này vẫn còn đặt ra nhiều
bài toán mà lời giải dường như thiết yếu cho sự phát triển tiếp theo của lý
thuyết.

Trong các ứng dụng, việc kiểm tra điều kiện (i) luôn tầm thường. Việc kiểm tra
(ii) tách thành hai phần: trường hợp $A'$ là một $A$-môđun tự do tương ứng
với kỹ thuật hạ xuống theo các cấu xạ phẳng ([3], các Định lý 1, 2 và 3) và
không gây khó khăn. Để xử lý trường hợp (b), ta sẽ dùng kết quả sau:

\result{ĐỊNH LÝ}{2} Cho $A$ là một vành địa phương Artin có iđêan tối đại
$\mathfrak m$, và $A'$ là một $A$-đại số chứa $A$, sao cho
$\mathfrak mA'\subset A$ và
\[
  A\longrightarrow A'\rightrightarrows A'\otimes_A A'
\]
là khớp (điều này đặc biệt xảy ra nếu $A'/A$ là một $A$-môđun có độ dài 1).
Cho $\mathcal F$ là phạm trù sợi ([3], A, Định nghĩa 1.1) của các bó tựa kết
hợp và phẳng trên các tiền lược đồ thay đổi. Khi đó cấu xạ
$\operatorname{Spec}(A')\longrightarrow\operatorname{Spec}(A)$ là một cấu
xạ $\mathcal F$-hạ xuống nghiêm ngặt ([3], A, Định nghĩa 1.7).

Nói cách khác, dữ liệu của một $A$-môđun phẳng $M$ hoàn toàn tương đương với
dữ liệu của một $A'$-môđun phẳng $M'$, được trang bị một
$A'\otimes_AA'$-đẳng cấu từ $M'\otimes_AA'$ lên $A'\otimes_AM'$ thỏa mãn
điều kiện bắc cầu thông thường của dữ liệu hạ xuống (loco citato).

Định lý 2 được chứng minh bằng cách trước hết chứng minh
\[
  H^i(A'/A,G_a)=0\qquad\text{với }i\geq 1
\]
([3], A, 4(e)), vì giả thiết $\mathfrak mA'\subset A$ cho phép dễ dàng quy về
trường hợp $A$ là một trường $k$ (cụ thể là $A/\mathfrak m$). Khi đó áp dụng
các tương đương được nêu trong [3], trang 16.

\fsection{C}{Ứng dụng vào một số trường hợp riêng}

\subsection*{1. Những nhận xét tổng quát về các hàm tử được biểu diễn bởi các tiền lược đồ}
\addcontentsline{toc}{subsection}{1. Những nhận xét tổng quát về các hàm tử được biểu diễn bởi các tiền lược đồ}

Cho $S$ là một tiền lược đồ Noether địa phương. Một tiền lược đồ $X$ trên
$S$ được gọi là \src{379}kiểu hữu hạn địa phương trên $S$ nếu, với mọi
$x\in X$ chiếu xuống $y\in Y$, tồn tại một lân cận afin của $y$, có vành $A$,
và một lân cận afin của $x$ nằm trên lân cận ấy, có vành $B$, sao cho $B$ là
một $A$-đại số kiểu hữu hạn. Có nhiều ví dụ quan trọng về các tiền lược đồ kiểu
hữu hạn địa phương trên $S$ nhưng không phải kiểu hữu hạn trên $S$, xuất hiện
như nghiệm của các bài toán phổ quát cổ điển; vì vậy cần có khả năng xem lược
đồ Picard của một đường cong như hợp của vô hạn thành phần liên thông (không
được nhầm chúng với thành phần liên thông của phần tử đơn vị, tức ``đa tạp
Picard''). Do đó đôi khi thuận tiện khi làm việc trong phạm trù $\mathcal C$
các tiền lược đồ kiểu hữu hạn địa phương trên $S$, để xét trong đó vấn đề biểu
diễn được của một hàm tử phản biến $F$. Mục tiêu chính của các bài trình bày
này là phát triển một kỹ thuật tổng quát cho phép nhận biết một hàm tử $F$ như
vậy có biểu diễn được hay không và nghiên cứu các tính chất của $S$-tiền lược
đồ $X$ tương ứng từ các tính chất của $F$. Nhân đây, trong nghiên cứu ấy xuất
hiện những ví dụ không bệnh lý về các tiền lược đồ trên $S$ không tách trên
$S$, đặc biệt dưới dạng các ``tiền lược đồ Picard'' của những $S$-lược đồ rất
tốt; vì thế không nên loại khỏi hình học đại số các tiền lược đồ không phải là
lược đồ.

Cho $X$ là một tiền lược đồ kiểu hữu hạn địa phương trên $S$, và đặt
\[
  F(Y)=\operatorname{Hom}_S(Y,X)
\]
là hàm tử phản biến liên kết. Ta có thể xét hạn chế $F_0$ của $F$ lên phạm trù
con $\mathcal C_0$ của $\mathcal C$ gồm các tiền lược đồ $Y$ trên $S$ vừa
Artin vừa hữu hạn trên $S$: khi $S=\operatorname{Spec}(\Lambda)$,
$\mathcal C_0$ chính là phạm trù đối ngẫu của phạm trù các $\Lambda$-đại số
Artin đã xét trong B. Nếu $Y=\operatorname{Spec}(A)$, với $A$ là một vành
Artin địa phương, thì $Y$ rút về một điểm duy nhất $y$ nằm trên một điểm đóng
$s$ của $S$, và một $S$-đồng cấu từ $Y$ đến $X$ (tức một phần tử của $F(Y)$)
được xác định bởi dữ liệu một điểm $x\in X$ nằm trên $s$ và một
$\mathcal O_{s,S}$-đồng cấu từ $\mathcal O_{x,X}$ đến $A$. Nếu tồn tại một
đồng cấu như vậy thì $x$ tất yếu là một điểm đóng của $X$ (vì trường thặng dư
của nó là đại số trên trường thặng dư của $s$). Điều này cho thấy hạn chế
$F_0$ của $F$ lên các ``$Y$-đại số Artin'' là pro-biểu diễn được, và được biểu
diễn bởi đại số tôpô trên $Y$ mà các thành phần địa phương là các đầy đủ hóa
$\widehat{\mathcal O}_{x,X}$ của các vành địa phương của $X$ tại những điểm
$x$ của $X$ vừa đóng vừa nằm trên các điểm đóng của $Y$. Vì thế, chỉ riêng
việc biết $F_0$ đã cung cấp thông tin chính xác về cấu trúc của $X$ (cụ thể là
cấu trúc các đầy đủ hóa của các vành địa phương tại các điểm vừa nêu). Cần lưu
ý rằng, ngay cả khi $S$

\src{380}
là phổ của một trường đóng đại số, chỉ bằng cách xét có hệ thống các ``đa
tạp'' $Y$ mà $\mathcal O_Y$ có thể chứa phần tử lũy linh, đặc biệt bằng cách
làm việc với phổ của các vành Artin địa phương, ta mới đạt được ``phát biểu
đúng'' của các bài toán phổ quát cổ điển và nắm bắt được phương diện ``vô cùng
bé'' của chúng.

Khi xuất phát từ một hàm tử $F$ cho trước và muốn quyết định nó có biểu diễn
được hay không, việc nghiên cứu hàm tử $F_0$ (nhờ các Định lý 1 và 2) sẽ cho
những chỉ dẫn gần như đầy đủ: hoặc, như thường xảy ra (chẳng hạn chỉ cần kiểm
tra bản chất của các tập $F(I_K,\xi)$ và cách chúng ứng xử theo hàm tử, xem A),
ta nhận thấy ngay $F_0$ không pro-biểu diễn được (điều này giải thích sự thất
bại của các cố gắng cho đến nay nhằm định nghĩa một cách ít nhiều tự nhiên
các đa tạp môđuli để phân loại các phân thớ vectơ hạng $>1$); hoặc ta kiểm tra
được $F_0$ biểu diễn được nhưng các không gian vectơ $F(I_K,\xi)$ không hữu
hạn chiều, khi đó phải bằng lòng với nghiệm ``hình thức''; hoặc cuối cùng,
$F_0$ được biểu diễn bởi một tích các vành địa phương đầy đủ Noether, điều này
cho những chỉ dấu rất mạnh rằng chính $F$ biểu diễn được và, kết hợp với các
tính chất tương tự nhưng mang tính toàn cục hơn mà chúng tôi dự định phát
triển sau này, chắc hẳn sẽ đủ để kéo theo điều đó. Cuối cùng, còn có những bài
toán hình học thú vị (xem các mục 4 và 5 dưới đây) trong đó ta chỉ có hàm tử
$F_0$ (không phát sinh từ một hàm tử ``toàn cục'' $F$ nào), và ta có lý do để
hài lòng nếu gắn được cho nó một ``lược đồ hình thức môđuli''.

Để kết thúc các nhận xét tổng quát này, hãy chỉ ra cách lý thuyết lược đồ giải
thích những hiện tượng có vẻ dị thường, chẳng hạn mặt Igusa $V$, có ``đa tạp
Picard'' $P$ rút về một điểm nhưng vẫn thỏa
\[
  H^1(V,\mathcal O_V)\ne 0 ;
\]
trong trường hợp này, $P$ là một nhóm ``thuần túy vô cùng bé'' không rút về
nhóm đơn vị, tức là nó được xác định bởi một đại số địa phương $O$ có hạng hữu
hạn trên trường cơ sở $k$ và được trang bị một ánh xạ đường chéo tương ứng với
cấu trúc nhân của $P$; nếu $\mathfrak m$ là iđêan cực đại của $O$, đối ngẫu của
$\mathfrak m/\mathfrak m^2$ đẳng cấu chính tắc với
$H^1(V,\mathcal O_V)$ (xem mục 3 dưới đây). Chỉ khi nhóm Picard là một nhóm
đại số theo nghĩa cổ điển, tức là trơn trên trường cơ sở $k$, thì chiều của
$H^1(V,\mathcal O_V)$ (luôn bằng chiều của $\mathfrak m/\mathfrak m^2$) mới
bằng chiều của nhóm Picard.

\subsection*{2. Các lược đồ $\underline{\operatorname{Hom}}_S(X,Y)$,
$\Pi_{X/S}Z$, $\underline{\operatorname{Aut}}_S(X)$, v.v.}
\addcontentsline{toc}{subsection}{2. Các lược đồ Hom, Pi và Aut}

Cho $X,Y$ là hai tiền lược đồ trên $S$; với mọi tiền lược đồ $T$ trên $S$,

\src{381}
đặt $X_T=X\times_ST$, $Y_T=Y\times_ST$, và xét tập
\[
  F(T)=\operatorname{Hom}_T(X_T,Y_T)
      =\operatorname{Hom}_S(X_T,Y)
      =\operatorname{Hom}_S(X\times_ST,Y)
\]
như một hàm tử phản biến theo $T$. Nếu hàm tử này biểu diễn được, ta ký hiệu
bởi $\underline{\operatorname{Hom}}_S(X,Y)$ tiền lược đồ trên $S$ biểu diễn
nó; do đó có một đẳng cấu hàm tử
\[
  \operatorname{Hom}_S
  \bigl(T,\underline{\operatorname{Hom}}_S(X,Y)\bigr)
  \xrightarrow{\sim}
  \operatorname{Hom}_S(T\times_SX,Y).
\]
Có nhiều biến thể của bài toán phổ quát này mà nghiệm quy về bài toán trên:
tiền lược đồ các $S$-tự đẳng cấu của một $S$-tiền lược đồ $X$ (đó sẽ là một
tiền lược đồ nhóm), tiền lược đồ các $S$-đồng cấu từ một $S$-tiền lược đồ nhóm
đến một tiền lược đồ nhóm khác (đó sẽ là một tiền lược đồ nhóm giao hoán nếu
lược đồ nhóm thứ hai giao hoán), v.v. Ta cũng có thể tổng quát hóa định nghĩa
của $\underline{\operatorname{Hom}}_S(X,Y)$ bằng cách xét một tiền lược đồ $Z$
trên tiền lược đồ $X$ trên $S$, cùng hàm tử
\[
  F(T)=\operatorname{Hom}_{X_T}(X_T,Z_T)
\]
(tập các ``tiết diện'' của $Z_T$ trên $X_T$); nếu hàm tử này biểu diễn được,
$S$-tiền lược đồ biểu diễn nó sẽ được ký hiệu $\Pi_{X/S}Z$; theo định nghĩa ta
có một đẳng cấu hàm tử
\[
  \operatorname{Hom}_S(T,\Pi_{X/S}Z)
  =\operatorname{Hom}_{X_T}(X_T,Z_T).
\]
Lấy $Z=Y\times_SX$, ta thu lại
$\underline{\operatorname{Hom}}_S(X,Y)$. Từ các định nghĩa này suy ra một bảng
công thức vừa hiển nhiên vừa hữu ích cho các tiền lược đồ mới đưa vào, mà ta
không trình bày ở đây (vì nó đúng trong mọi phạm trù có các tích và tích thớ).
Nghiêm túc hơn là câu hỏi về sự tồn tại của các lược đồ dạng
$\underline{\operatorname{Hom}}_S(X,Y)$. Trước hết, ta thấy rằng, với $X$ cố
định, $\underline{\operatorname{Hom}}_S(X,Y)$ chỉ có thể tồn tại với mọi $Y$
trên $S$ (hoặc $\Pi_{X/S}Z$ với mọi $Z$ trên $X$) nếu $X$ phẳng trên $S$.
Hơn nữa, chỉ hợp lý khi kỳ vọng có nghiệm với $Y$ đủ tổng quát nếu $X$ còn
riêng trên $S$. Mặt khác, dường như các điều kiện này đủ cho sự tồn tại của
$\underline{\operatorname{Hom}}_S(X,Y)$ và $\Pi_{X/S}Z$, miễn là áp đặt, khi
cần, một giả thiết kiểu ``tựa xạ ảnh'' lên $Y/S$, tương ứng lên $Z/X$; ít nhất
điều này có thể kiểm tra trong khá nhiều tình huống (chẳng hạn khi $Y$ afin
trên $S$; hoặc, bằng các phép dựng trực tiếp sơ cấp, khi $X$ hữu hạn trên $S$).
Sau đây là hệ quả của các Định lý 1 và 2:

\src{382}
\result{MỆNH ĐỀ}{2.1} Cho $\Lambda$ là một vành Noether, $X$ và $Y$ là hai
tiền lược đồ tùy ý trên $\Lambda$, và xét hàm tử
\[
  F(A)=\operatorname{Hom}_A(X_A,Y_A)
\]
trên phạm trù $\mathcal C_0$ các $\Lambda$-đại số Artin. Nếu $X$ phẳng trên
$\Lambda$, hàm tử này pro-biểu diễn được.

Hơn nữa, ta kiểm tra được rằng với mọi $A\in\mathcal C_0$ và $\xi\in F(A)$ có
một đẳng cấu chính tắc
\[
  F(I_A,\xi)=H^0\!\left(
    X_A,
    \underline{\operatorname{Hom}}_{\mathcal O_{X_A}}
    \bigl(\xi^*(\Omega^1_{Y_A/A}),\mathcal O_{X_A}\bigr)
  \right),
\]
trong đó $\Omega^1_{Y_A/A}$ là bó các $1$-vi phân Kähler của $Y_A$ đối với
$A$. Lấy $A$ là một trường rồi dùng A.5.1 và định lý hữu hạn của [2], ta được
hệ quả sau:

\result{HỆ QUẢ}{} Giả sử $X$ phẳng và riêng trên $S$, còn $Y$ kiểu hữu hạn địa
phương trên $S$. Khi đó $F$ pro-biểu diễn được và các thành phần địa phương của
đại số tôpô tương ứng trên $\Lambda$ là các vành Noether.

\result{NHẬN XÉT}{} Các bài toán được xét trong mục này, cùng nhiều bài toán
khác, theo truyền thống được nghiên cứu trong khuôn khổ hình học đại số cổ điển
bằng ``tọa độ Chow'' của các chu trình trong không gian xạ ảnh, nhờ đó có thể
xem các chu trình ấy như những điểm của các đa tạp xạ ảnh thích hợp. Thủ tục
này, và nói chung việc dùng tọa độ Chow, dường như không thể tránh khỏi sự bất
đầy đủ theo quan điểm lược đồ, vì nó tiêu hủy các phần tử lũy linh trong các đa
tạp tham số và đặc biệt không thích hợp cho việc nghiên cứu thỏa đáng các biến
phân vô cùng bé của chu trình (chưa kể tính không nội tại do phụ thuộc vào
không gian xạ ảnh). Đáng tiếc, ngôn ngữ tọa độ Chow từng là ngôn ngữ duy nhất
được nhiều nhà hình học đại số dùng để nghiên cứu các họ đa tạp hoặc các họ
chu trình, và điều đó dường như đã tạo ra một trở ngại nghiêm trọng cho việc
làm sáng tỏ những khái niệm này, bất chấp lợi ích kỹ thuật không thể phủ nhận
(có lẽ chỉ tạm thời) của chúng. Nếu muốn có cái tương tự các đa tạp Chow trong
lý thuyết lược đồ, ta đi đến bài toán phổ quát sau: cho $X$ là một tiền lược đồ
trên $S$; với mọi tiền lược đồ $T$ trên $S$, xét tập $F(T)$ các tiền lược đồ
con đóng của $X_T=X\times_ST$ phẳng trên $T$. Ta tìm cách biểu diễn hàm tử theo
$T$ này bằng một tiền lược đồ trên $S$. Tổng quát hơn, có thể cho một bó tựa
kết hợp $G$ trên $X$ và lấy $F(T)$ là tập các bó

\src{383}
thương của $G_T$ phẳng trên $T$. Dường như bài toán có nghiệm bằng một lược đồ
$C$ kiểu hữu hạn địa phương trên $S$ khi $X$ riêng trên tiền lược đồ Noether
địa phương $S$ và $G$ còn coherent. Dù sao, chỉ cần giả sử $S$ Noether địa
phương thì hạn chế của $F$ lên các ``$S$-đại số Artin'' là pro-biểu diễn được;
và nếu thêm rằng $X$ riêng trên $S$ và $G$ coherent thì các thành phần địa
phương của vành tôpô tương ứng $O$ là Noether. Dĩ nhiên, ngay cả sau khi chứng
minh sự tồn tại của ``lược đồ Chow'' của $X$ trên $S$, vẫn còn phải tìm một
phân tích thành các tập mở rời nhau $C_i$ (tương ứng với việc cố định các bất
biến liên tục như bậc và chiều của các chu trình đang biến thiên) kiểu hữu hạn
(chứ không chỉ kiểu hữu hạn địa phương) trên $S$, làm rõ quan hệ giữa lược đồ
này với các đa tạp Chow cổ điển, và xác định khi nào một $C_i$ thậm chí xạ ảnh,
hoặc ít nhất tựa xạ ảnh, trên $S$.

\subsection*{3. Các lược đồ Picard}
\addcontentsline{toc}{subsection}{3. Các lược đồ Picard}

Cho $f:X\longrightarrow S$ là một $S$-tiền lược đồ, xét bó nhân
$\mathcal O_X^*$ các đơn vị của bó cấu trúc của $X$, và nhóm
\[
  P(X/S)=H^0\bigl(S,R^1f_*(\mathcal O_X^*)\bigr),
\]
gọi là nhóm Picard tương đối của $X/S$. Như vậy, một phần tử của nhóm này được
xác định bằng cách cho một phủ mở $(U_i)$ của $S$ và một bó khả nghịch $L_i$
trên mỗi $f^{-1}(U_i)$, sao cho với mọi cặp $(i,j)$,
$L_i|f^{-1}(U_i\cap U_j)$ đẳng cấu với
$L_j|f^{-1}(U_i\cap U_j)$, ít nhất là địa phương trên $U_i\cap U_j$ (tức hai bó
này ``tương đương'' theo nghĩa của [3], B, 4). Nếu $X/S$ có một tiết diện thì
$P(X/S)$ không gì khác hơn tập các lớp bó khả nghịch trên $X/S$ modulo
``tương đương'' (loco citato). Bây giờ, với mọi $T$ trên $S$, đặt
\[
  F(T)=P(X_T/T) ;
\]
ta được một hàm tử hiệp biến theo $T$, có thể gọi là hàm tử Picard của $X/S$;
nếu hàm tử này biểu diễn được thì tiền lược đồ trên $S$ biểu diễn nó được gọi
là tiền lược đồ Picard của $X/S$ và ký hiệu $\underline P(X/S)$. Do đó có một
đẳng cấu hàm tử
\[
  \operatorname{Hom}_S\bigl(T,\underline P(X/S)\bigr)
  \xrightarrow{\sim}P(X_T/T).
\]
Phép lập tiền lược đồ Picard tương thích với mở rộng cơ sở; đặc biệt, các tiền
lược đồ Picard của các thớ của $X$ trên $S$ (là các tiền lược đồ trên các
trường thặng dư $K(s)$ của $s\in S$) chính là các thớ của
$\underline P(X/S)$. Dĩ nhiên, vì $P(X_T/T)=F(T)$ là một nhóm giao hoán, các
tiền lược đồ Picard là các tiền lược đồ nhóm. Cũng cần lưu ý rằng các Jacobian
tổng quát của Rosenlicht không gì khác hơn các thành phần liên thông của đơn
vị trong các lược đồ Picard của những đường cong đầy đủ có thể có kỳ dị; điều
này hẳn làm hiển nhiên phần lớn các tính chất của chúng (một khi đã chứng minh
sự tồn tại).

\result{NHẬN XÉT}{} Định nghĩa được chọn ở đây chỉ hợp lý khi mọi điểm của
$Y$ có một lân cận mở $U$ mà trên đó $X$ có một tiết diện. Trong trường hợp
tổng quát, phải sửa nhẹ định nghĩa hàm tử Picard để vẫn có thể thu được một
định lý tồn tại.

Ở đây, các điều kiện hợp lý cho sự tồn tại của một tiền lược đồ Picard là:
$X$ riêng và phẳng trên $S$, $f_*(\mathcal O_X)=\mathcal O_S$, và $X$ có địa
phương một tiết diện trên $S$. Điều kiện này xuất hiện tự nhiên khi áp dụng kỹ
thuật hạ xuống: để loại bỏ các tự đẳng cấu của một bó khả nghịch $L$ trên $X$,
ta trang bị cho nó một tiết diện được đánh dấu trên tiết diện $s$ ([3], B, 4).
Đặc biệt, ta thu được:

\src{384}
\result{MỆNH ĐỀ}{3.1} Giả sử $X$ phẳng trên
$S=\operatorname{Spec}(\Lambda)$, với $\Lambda$ Noether, và với mọi $T$ kiểu
hữu hạn trên $S$ ta có
\[
  f_{T*}(\mathcal O_{X_T})=\mathcal O_T
\]
(nếu $f$ riêng và có các thớ khả tách, hoặc nếu $S$ là phổ của một trường,
thì theo Künneth điều kiện cuối tương đương với
$f_*(\mathcal O_X)=\mathcal O_S$). Khi đó hàm tử Picard của $X/S$ trên phạm
trù các $\Lambda$-đại số Artin là pro-biểu diễn được.

Hơn nữa, ở đây ta có
\[
  F(I_A,\xi)=H^1(X_A,\mathcal O_{X_A}),
\]
đặc biệt:

\result{HỆ QUẢ}{} Nếu $X$ riêng trên $S$, thì các thành phần địa phương của
$\Lambda$-đại số tôpô $O$ tương ứng với hàm tử Picard là Noether.

\result{NHẬN XÉT}{} Ta có thể tổng quát hóa các định nghĩa và kết quả của mục
này cho việc phân loại các phân thớ chính trên $X$, với nhóm cấu trúc là một
lược đồ nhóm $G$ trên $S$ vừa afin, vừa phẳng trên $S$, vừa giao hoán. Khi $G$
không giao hoán, phân thớ nhóm liên hợp của một phân thớ chính (mà các tiết
diện là các tự đẳng cấu của phân thớ chính) không còn tầm thường, nên Mệnh đề
3.1 không còn đúng nguyên dạng. Tuy nhiên, ta có thể sửa bài toán phổ quát sao
cho nó vẫn có nghiệm (ít nhất, hiện tại, trong hình học hình thức). Quy tắc
vàng cần ghi nhớ, trong bối cảnh của \src{385}mục này và các mục sau, mỗi khi
ta tìm các ``lược đồ môđuli'' cho những lớp đối tượng chỉ được xác định sai
khác đẳng cấu, luôn là: loại bỏ các tự đẳng cấu khả dĩ của các đối tượng cần
phân loại bằng cách đưa vào, nếu cần, những cấu trúc phụ (các điểm hay phần tử
của tiết diện được đánh dấu, cố định các dạng vi phân, v.v.), được chọn đủ vô
hại để không làm thay đổi đáng kể bài toán ban đầu.

\subsection*{4. Môđuli hình thức của một đa tạp}
\addcontentsline{toc}{subsection}{4. Môđuli hình thức của một đa tạp}

Cho $\Lambda$ là một vành Noether địa phương có trường thặng dư $k$ (thông
thường nhất, $\Lambda$ bằng $k$, hoặc bằng một $p$-vành Cohen), và cho $X_0$
là một tiền lược đồ trên $k$. Với mọi $\Lambda$-đại số Artin địa phương $A$,
xét tập $F(A)$ các lớp, sai khác đẳng cấu, của các $A$-tiền lược đồ $X$ phẳng
trên $A$, được trang bị một đẳng cấu
\[
  (*)\qquad
  X\otimes_A\kappa(A)\xleftarrow{\sim}
  X_0\otimes_k\kappa(A),
\]
trong đó $\kappa(A)$ là trường thặng dư của $A$; dĩ nhiên, các đẳng cấu giữa
những $A$-tiền lược đồ phẳng như vậy phải tôn trọng đẳng cấu vừa nêu, vốn là
một phần của cấu trúc. Nếu $A$ là một $\Lambda$-đại số Artin không nhất thiết
địa phương, có các thành phần địa phương $A_i$, ta lấy $F(A)$ là tích các
$F(A_i)$. Như vậy, $F$ trở thành một hàm tử nhân tính theo $A$, có thể gọi là
hàm tử môđuli của $X_0$ (và $\Lambda$). Nếu hàm tử này biểu diễn được, nó tương
ứng với một $\Lambda$-đại số tôpô địa phương $O$ có trường thặng dư $k$, và
phổ hình thức của $O$ được gọi là lược đồ hình thức môđuli của $X_0$ (và
$\Lambda$); xem [2] để biết một số chi tiết về khái niệm này.

Ở đây, để áp dụng kỹ thuật hạ xuống, các tự đẳng cấu ``hữu hạn'' của $X_0$ là
vô hại, vì chúng không ảnh hưởng đến sự tồn tại của các tự đẳng cấu (theo
nghĩa đã nêu trên) của các $A$-tiền lược đồ $X$; khi $A$ không rút về một
trường, điều kiện cần và đủ để $X$ không có tự đẳng cấu $A$ không tầm thường
nào là
\[
  H^0(X_0,\Theta_{X_0/k})=0,
\]
trong đó $\Theta_{X_0/k}$ ký hiệu bó các $k$-đạo hàm (= bó tiếp tuyến) của
$X_0$. Mặt khác, ta dễ dàng thu được (ít nhất nếu $X_0$ trơn trên $k$)
\[
  F(I_A,\xi)=H^1(X_A,\Theta_{X_A/A}).
\]
Như thường lệ, suy ra:

\src{386}
\result{MỆNH ĐỀ}{4.1} Giả sử
$H^0(X_0,\Theta_{X_0/k})=0$. Khi đó tồn tại lược đồ hình thức môđuli của
$X_0$. Nếu thêm rằng $X_0$ riêng trên $k$, lược đồ hình thức môđuli là
Noether.

\result{NHẬN XÉT}{} 1\textsuperscript{o} Khi không giả sử $X_0$ trơn trên
$k$, $F(I_A,\xi)$ được đồng nhất với một $A$-môđun con của
\[
  \operatorname{Ext}^1_{\mathcal O_{P_A}}
  \bigl(P_A;\mathcal I_{X_A},\mathcal O_{X_A}\bigr),
\]
trong đó đặt $P_A=X_A\times_\Lambda X_A$, xem $\mathcal O_{X_A}$ như một bó
coherent trên $P_A$ thông qua cấu xạ đường chéo
$X_A\longrightarrow P_A$, và $\mathcal I_{X_A}$ ký hiệu bó iđêan coherent
trên $P_A$ được xác định bởi cấu xạ đường chéo. Chính xác hơn, một phép toàn
cục hóa đơn giản của lý thuyết Hochschild cho thấy nhóm
$\operatorname{Ext}^1$ trên đồng nhất với tập các lớp, sai khác đẳng cấu, của
các bó $I_A$-đại số phẳng $\mathcal O$ trên $X_A$, được trang bị một đẳng cấu
tăng cường
\[
  \mathcal O\otimes_{I_A}A\longrightarrow\mathcal O_{X_A}
\]
(nhắc lại rằng $I_A=A[t]/(t^2)$). Môđun con $F(I_A,\xi)$ là môđun tương ứng
với các bó đại số giao hoán. Vì vậy, trái với điều [2] gợi ra, các giả thiết
trơn không thiết yếu trong lý thuyết môđuli.

2\textsuperscript{o} Nhắc lại (loco citato) rằng, đặc biệt, mọi đường cong đại
số $X_0$ trơn và riêng trên $k$ đều có một lược đồ hình thức môđuli trơn trên
$\Lambda$, có chiều tương đối $3g-3$ nếu giống $g$ là $\geq2$, và bằng $g$ nếu
$g=0$ hoặc $1$. Hai trường hợp sau không còn trực tiếp thuộc Mệnh đề 4.1. Tuy
nhiên, trong trường hợp các đường cong elliptic ($g=1$), có thể quy chúng về
mệnh đề ấy bằng các xét đoán sau đây.

Dĩ nhiên, ta có thể biến đổi Mệnh đề 4.1 tùy ý bằng cách xét các hệ lược đồ
trên $k$ được trang bị nhiều cấu trúc khác nhau. Chẳng hạn, giả sử $X_0$ là
một lược đồ Abel trên $k$, có điểm gốc được đánh dấu (tức là xem $X_0$ như một
lược đồ nhóm trên $k$), và cho $F(A)$ là tập các lớp, sai khác đẳng cấu, của
các lược đồ Abel trên $A$ (tức các lược đồ nhóm riêng và trơn trên $A$) được
trang bị một đẳng cấu lược đồ Abel $(*)$. Ta kiểm tra được rằng việc áp đặt một
cấu trúc nhân (thậm chí chỉ một ``tiết diện đơn vị'') loại bỏ các tự đẳng cấu
vô cùng bé, và do đó tồn tại một lược đồ hình thức môđuli tương ứng với một
vành địa phương đầy đủ Noether $O$. Mặt khác, có thể chứng minh rằng nếu $X$
là một lược đồ riêng và trơn, có các thớ ``liên thông tuyệt đối'' trên một tiền
lược đồ Noether địa phương $S$, thì mọi cấu trúc nhân trên $X$ có một tiết
diện đơn vị đều tất yếu kết hợp và giao hoán (miễn là nó như vậy trên ít nhất
một thớ và $S$ liên thông), và hơn nữa được xác định duy nhất bởi việc biết

\src{387}
tiết diện đơn vị. Ngoài ra, giả sử $S$ là phổ của một vành Artin địa phương
$A$ có trường thặng dư $k$, $X$ riêng trên $A$ và được trang bị một tiết diện
$s$, và cuối cùng $X\otimes_Ak$ được trang bị cấu trúc lược đồ Abel trên $k$
mà phần tử trung hòa là điểm của $X\otimes_Ak$ tương ứng với $s$; một phép
tính cản trở đơn giản, kết hợp với một lập luận của Serre, cho phép chứng minh
rằng thực sự tồn tại trên $X$ một cấu trúc nhân nhận $s$ làm tiết diện đơn vị.
(Từ đó, dùng ``định lý tồn tại'' của [2] để chuyển sang trường hợp $A$ Noether
địa phương đầy đủ, rồi dùng kỹ thuật hạ xuống của [3] cho trường hợp tổng quát,
ta có thể chứng minh phát biểu tương tự với mọi $S$ Noether địa phương liên
thông.) Điều này chứng minh rằng hàm tử $F(A)$ đang xét ở đây đẳng cấu với hàm
tử tương tự được định nghĩa ở đầu mục này khi bỏ qua cấu trúc nhân trên $X_0$.
Đặc biệt, suy ra rằng nếu $\mathfrak m$ là iđêan cực đại của $O$ thì
$\mathfrak m/\mathfrak m^2$ đẳng cấu chính tắc với đối ngẫu của
$H^1(X_0,\Theta_{X_0/k})$, và do đó có chiều $n^2$, trong đó
$n=\dim X_0$. Sẽ rất thú vị nếu xác định được $O$ có trơn trên $\Lambda$ hay
không, tức có đẳng cấu với một đại số chuỗi hình thức theo $n^2$ biến trên
$\Lambda$ hay không. Mục 1, Mệnh đề 5.2, cho phép phát biểu tương đương bài
toán này như một bài toán tồn tại các lược đồ Abel rút về một lược đồ Abel cho
trước. Dù sao, các phương pháp siêu việt cho thấy câu trả lời là khẳng định nếu
$k$ có đặc số không. Trong đặc số $p\ne0$, hiển nhiên chỉ cần xét trường hợp
$\Lambda$ là vành vectơ Witt xây dựng trên một trường đóng đại số $k$. Có lẽ
đã đến lúc ``hàm tử Greenberg'' chứng tỏ giá trị của mình\ldots

\subsection*{5. Mở rộng các phủ}
\addcontentsline{toc}{subsection}{5. Mở rộng các phủ}

Cho $\mathfrak X$ là một tiền lược đồ hình thức Noether [2] và $U$ là một
phần mở của $\mathfrak X$, được xác định địa phương bởi sự ``không triệt tiêu''
của một tiết diện của $\mathcal O_{\mathfrak X}$ không là ước của không trên
$\mathfrak X$, và vì thế đủ lớn để mọi tiết diện của
$\mathcal O_{\mathfrak X}$ trên một tập mở $V$ triệt tiêu trên $U\cap V$ đều
triệt tiêu. Cho $\mathcal J$ là một ``iđêan xác định'' của $\mathfrak X$, và
đặt
\[
  X_n=(\mathfrak X,\mathcal O_{\mathfrak X}/\mathcal J^{n+1}),
\]
đây là một tiền lược đồ Noether thông thường. Nếu
$\mathfrak X',\mathfrak X''$ là hai phủ phẳng của $\mathfrak X$ (tức hai tiền
lược đồ trên $\mathfrak X$ được xác định bởi các bó đại số coherent và tự do
địa phương khi xem như các bó môđun), không phân nhánh trên $U$, thì ánh xạ
hiển nhiên
\[
  \operatorname{Hom}_{\mathfrak X}(\mathfrak X',\mathfrak X'')
  \longrightarrow
  \operatorname{Hom}_{X_0}(X'_0,X''_0)
\]
là đơn ánh; đặc biệt, một tự đẳng cấu của $\mathfrak X'$ cảm sinh đồng nhất
trên $X'_0$ chính là đồng nhất. Điều này cho phép áp dụng kỹ thuật hạ xuống vào
tình huống này. Đặc biệt, hãy xuất phát từ một phủ phẳng $X'_0$ của $X_0$,
không phân nhánh

\src{388}
trên $U_0$; ký hiệu $G(\mathfrak X)$ là tập các lớp, sai khác đẳng cấu (các
đẳng cấu cảm sinh đồng nhất trên $X'_0$), của các phủ phẳng $\mathfrak X'$ của
$\mathfrak X$ cảm sinh $X'_0$ trên $X_0$ (và do đó tất yếu không phân nhánh
trên $U$). Định nghĩa tương tự $G(V)$ với mọi phần mở $V$ của $\mathfrak X$,
và tổng quát hơn $G(\mathfrak Y)$ với mọi tiền lược đồ hình thức
$\mathfrak Y$ trên $\mathfrak X$. Khi ấy, các kết quả của [2] và [3] trước hết
kéo theo những kết quả sau:

\noindent a. Khi $V$ chạy qua các tập mở của $\mathfrak X$, các $G(V)$ lập
thành một bó trên $\mathfrak X$, ký hiệu $G_{\mathfrak X}=G$. Hạn chế của bó
này lên $U$ là bó hằng có các thớ rút về một phần tử.

Tổng quát hơn, việc xác định các thớ của $G_{\mathfrak X}$ là một vấn đề về
các vành địa phương đầy đủ; chính xác hơn:

\noindent b. Với mọi $x\in\mathfrak X$, ta có
\[
  G_x=G\bigl(\operatorname{Spec}(\widehat{\mathcal O}_{\mathfrak X,x})\bigr),
\]
(= các lớp, sai khác đẳng cấu, của các đại số $B$ hữu hạn và tự do trên
$\widehat{\mathcal O}_{\mathfrak X,x}$, được trang bị một đẳng cấu của
$B\otimes_{\widehat{\mathcal O}_{\mathfrak X,x}}\mathcal O_{X_0,x}$ với
$\mathcal O'_{0x}$, trong đó $\mathcal O'_0$ là bó đại số trên $X_0$ xác định
$X'_0$).

\noindent c. Có một đẳng cấu chính tắc
\[
  G_{\mathfrak X}=\varprojlim_nG_{X_n} ;
\]
nói cách khác, với mọi tập mở $V$ của $\mathfrak X$, ta có
\[
  G(V)=\varprojlim_nG(V_n).
\]

\noindent d. Giả sử $\mathfrak X$ thu được từ một lược đồ thông thường $X$
riêng trên một vành Noether địa phương đầy đủ $\Lambda$ có iđêan xác định
$\mathfrak m$, bằng cách lấy đầy đủ hóa $\mathcal J$-adic của $\mathcal O_X$,
trong đó $\mathcal J=\mathfrak m\mathcal O_X$. Khi đó $G(\mathfrak X)$ đẳng
cấu chính tắc với tập các lớp phủ phẳng của lược đồ thông thường $X$ ``rút về
$X'_0$''.

Nói một cách hình tượng, có thể nói rằng (a) và (b) thiết lập các quan hệ cơ
bản giữa phương diện địa phương và phương diện toàn cục của bài toán, (c) cho
các quan hệ giữa phương diện ``hữu hạn'' và phương diện ``vô cùng bé'', còn
(d) nhắc lại (dưới các điều kiện đã nêu) sự đồng nhất giữa phương diện ``hình
thức'' và phương diện ``đại số''.

Bây giờ giả sử $\mathfrak X$ được xác định trên một vành Noether địa phương
đầy đủ $\Lambda$, với $\mathcal J=\mathfrak m\mathcal O_{\mathfrak X}$ (do đó
$X_0$ là một tiền lược đồ trên $\Lambda/\mathfrak m$). Với mọi đại số $A$ hữu
hạn trên $\Lambda$, đặt
\[
  F(A)=G(\mathfrak X\times_\Lambda A).
\]
Đây là một hàm tử hiệp biến theo $A$, nhận giá trị trong phạm trù các tập hợp,
và theo (c) hàm tử này được xác định hoàn toàn khi biết nó trên các $A$ Artin;
nói rằng hàm tử này pro-biểu diễn được, tức có dạng
\[
  F(A)=\operatorname{Hom}_{\Lambda\text{-algèbres top.}}(O,A),
\]

\src{389}
trong đó $O$ là một đại số tôpô trên $\Lambda$ thuộc loại đã xét ở A, mục 5,
tương đương với việc điều ấy đúng khi hạn chế vào các $A$ Artin. Khi đó, sự kết
hợp của các Định lý 1 và 2 thực sự kéo theo:

\result{MỆNH ĐỀ}{5.1} Hàm tử trên pro-biểu diễn được.

Dĩ nhiên, theo (a), nếu $U=\mathfrak X$, $G(\mathfrak Y)$ rút về một phần tử
duy nhất với mọi $\mathfrak Y$ trên $\mathfrak X$, và khi đó hàm tử $F$ ít
đáng quan tâm (ta có $O=\Lambda$). Dường như trong hầu hết mọi trường hợp khác,
vành địa phương tôpô $O$ không Noether. Tuy nhiên, sự tồn tại của nó làm nổi
bật một cách đáng ngạc nhiên bản chất ``liên tục'' của tập nghiệm
$G(\mathfrak X)$ (tương ứng trực giác với việc có một lựa chọn ``liên tục''
trong cách sự phân nhánh lan truyền khi mở rộng $X'_0$). Hãy so sánh kết quả
này với quan điểm của J.-P. Serre [4] trong lý thuyết trường lớp địa phương,
vốn cũng làm nổi bật tính liên tục của nhóm Galois tôpô của mở rộng Abel cực
đại của một trường địa phương ``hình học''; ở đó nhóm đối ngẫu (theo nghĩa
Pontryagin) xuất hiện như một giới hạn quy nạp của các nhóm đại số (hoặc ít
nhất tựa đại số); ở đây cũng vậy, việc phân loại các mở rộng được cho bởi các
``đa tạp'' vô hạn chiều. Mặt khác, trong phần trên ta có thể lấy
$\mathfrak X$ là phổ hình thức của một vành địa phương đầy đủ (mà $\Lambda$
chẳng hạn là một vành con Cohen), và có thể hy vọng các kết quả của mục này
được dùng để nghiên cứu các mở rộng của một vành địa phương đầy đủ có chiều
$>1$. Cả trong trường hợp địa phương lẫn toàn cục, có lẽ chúng cho phép phát
biểu các quan hệ chính xác giữa những hiện tượng phân nhánh cao và những hiện
tượng ở đặc số không (có thể tiếp cận bằng phương pháp siêu việt). Dù sao,
chính phân tích sơ bộ của Mệnh đề 5.1 cho phép mở rộng các phương pháp trình
bày trong [2] để nghiên cứu nhóm cơ bản sang trường hợp ``phân nhánh vừa'' và
giải ``bài toán ba điểm'' bằng phương pháp siêu việt.

Để kết thúc, hãy lưu ý rằng tình huống trở nên đơn giản khi $X_0$ có chiều 1:
khi ấy, theo (a) và (b),
\[
  G(\mathfrak X)\simeq
  \prod_iG\bigl(\operatorname{Spec}(\widehat{\mathcal O}_{\mathfrak X,x_i})\bigr),
\]
trong đó $x_i$ là các điểm của $X_0-U$: ta có thể cho tùy ý các mở rộng ``địa
phương'' tại những điểm phân nhánh. Ngoài ra, khi $X_0$ chuẩn tắc, ta nhận thấy
lược đồ hình thức môđuli được bảo đảm bởi 5.1 là trơn trên
$\operatorname{Spec}(\Lambda)$.

\src{390}
\begin{center}
  \textbf{TÀI LIỆU THAM KHẢO}
\end{center}

\begin{description}
  \item[{[1]}] Grothendieck (Alexander). --- \emph{Sur quelques points
    d'algèbre homologique}, Tohoku Math. J., t. 9, 1957, tr. 119--221.
  \item[{[2]}] Grothendieck (Alexander). --- \emph{Géométrie formelle et
    géométrie algébrique}, Séminaire Bourbaki, t. 11, 1958/59, fasc. 3,
    số 182.
  \item[{[3]}] Grothendieck (Alexander). --- \emph{Technique de descente et
    théorèmes d'existence en géométrie algébrique, I : Généralités, Descente
    par morphismes fidèlement plats}, Séminaire Bourbaki, t. 12, 1959/60,
    fasc. 1, số 190.
  \item[{[4]}] Serre (Jean-Pierre). --- \emph{Corps locaux et isogénies},
    Séminaire Bourbaki, t. 11, 1958/59, fasc. 3, số 185.
\end{description}

\par
